\section{The origin region}\label{sec:origin}

In this section we exclude a rightmost zero close to the upper endpoint of
the unit disk.  The estimates apply down to arbitrarily small positive
scales.

\begin{proposition}[Origin region]\label{or:theorem}
Let $P$ have degree $n\ge6251$, with all its zeros in $\overline\D$.
Suppose that $a>0$ and
\[
             |a-\xi|>1\qquad(P'(\xi)=0).
\]
Let $b=x+iy$ be a rightmost zero of $P$, where $x\ge a$ and $y\ge0$.
Then
\begin{equation}\label{or:scale-conclusion}
       \max(x,1-y)>\rho_*:=\frac{30501}{500000}.
\end{equation}
All zeros and critical points in this statement, and in the sums below,
are counted with their algebraic multiplicities.
\end{proposition}

\subsection{Normalization and radial identities}

Suppose, to the contrary, that
\[
 \rho=\max(x,1-y)\le\rho_*,\qquad
 F(u)=P(i+\rho u),\qquad
 \beta=X-iT=\frac{b-i}{\rho},\qquad \alpha=\frac a\rho.
\]
Then
\begin{equation}\label{or:normalization}
 0<\alpha\le X\le1,\qquad 0<T\le1,\qquad \max(X,T)=1.
\end{equation}
Here $T>0$: a point $x+iy\in\overline\D$ with $x>0$ cannot have
$y=1$.  Every zero $u$ of $F$ satisfies
\begin{equation}\label{or:original-domain}
                  \Re u\le X,\qquad |i+\rho u|\le1.
\end{equation}
Gauss--Lucas and the strict hole give, for each critical point $v=r-it$,
\begin{equation}\label{or:critical-domain}
 \frac\rho2|v|^2\le t<\frac\rho2|v-\alpha|^2,
                  \qquad r<\frac\alpha2.
\end{equation}
The root $\beta$ is simple, since its real part is strictly larger than
that of every critical point.  Put
\begin{equation}\label{or:notation}
 \begin{gathered}
 N=n-1,\quad N_0=6250,\quad n_0=6251,\quad c_\rho=\rho/2,\quad
 I(u)=-\Im u-\frac\rho2|u|^2,\\
 S=\frac{F''(\beta)}{F'(\beta)}=\sum_v\frac1{\beta-v},\quad
 h=|S|,\quad Q=\sum_v\frac1{|\beta-v|^2},\quad
 d=\min_v|\beta-v|,\\
 y_0=\frac Qh,\quad q=\frac Q{h^2},\quad
 \mu=\frac1{dh},\quad D=dh,\quad \tau=\frac{N\rho}{2}.
 \end{gathered}
\end{equation}
We use $y_0$ for the moment ratio, to distinguish it from the physical
coordinate $y$ of $b$.  The original and critical insets $I$ are
nonnegative.  Moreover,
\begin{equation}\label{or:horizontal}
 \Re S>0,\quad h>0,\quad
 Q<\frac{2\Re S}{X},\quad h^2\le NQ,\quad \frac hn<y_0.
\end{equation}
Indeed, $X-r>X/2$ in every summand of $\Re S$; the penultimate
inequality is Cauchy--Schwarz.  Also, by
\eqref{or:critical-domain},
\begin{equation}\label{or:inset-critical}
 I(v)<\rho\alpha(\alpha/2-r)\le\rho X(X/2-r).
\end{equation}
For the last comparison subtract the two quadratic expressions: their
difference is $(X-\alpha)((X+\alpha)/2-r)\ge0$.

For $b=i+\rho\beta$ and any $z\ne\beta$,
\begin{equation}\label{or:radial-one}
 \Re\frac b{\beta-z}
       =\frac\rho2+\frac{I(z)-I(\beta)}{|\beta-z|^2}.
\end{equation}
Summing over the critical points, and writing
$E_c=\sum_v I(v)/|\beta-v|^2$, gives
\begin{equation}\label{or:radial}
 \tau=\rho X\Re S-(1-\rho T)\Im S-E_c+I(\beta)Q.
\end{equation}
Since $E_c\ge0$, $|b|\le1$, and $I(\beta)<T$, it follows that
\begin{equation}\label{or:radial-unphased}
                         \tau<h(1+Ty_0).
\end{equation}
On the other hand \eqref{or:inset-critical} gives
\[
            E_c<\rho(X\Re S-X^2Q/2).
\]
Substitution into \eqref{or:radial}, retaining the term $-E_c$, yields
\begin{equation}\label{or:radial-sharp}
 T(1-\rho T/2)Q<\tau+(1-\rho T)\Im S.
\end{equation}
Thus, if $|\Im S|<J$,
\begin{align}
 \tau&<\rho h+J+Ty_0h,\label{or:radial-phase}\\
 \left(\frac hn\right)^2
 &<\left(\frac hN\right)^2
 <\frac{\rho/2+J/N}{T(1-\rho T/2)}.\label{or:radial-ratio}
\end{align}

At zero inset we have the following bound.
\begin{lemma}[A boundary source]\label{or:zero-inset}
If $X=1$ and $I(\beta)=0$, then $h>N/2$.
\end{lemma}
\begin{proof}
Now $|b|=1$.  At a simple zero the original-root factorization gives
$S=2\sum_{u\ne\beta}(\beta-u)^{-1}$.  If $z=i+\rho u$ is another
physical zero, then
\[
 2\Re\frac b{b-z}-1=\frac{1-|z|^2}{|b-z|^2}\ge0.
\]
Consequently
\[
 \Re(bS)=2\rho\sum_{z\ne b}\Re\frac b{b-z}\ge N\rho.
\]
The critical-root version of \eqref{or:radial-one} says
$\Re(bS)=N\rho/2+E_c$.  Hence
\[
 N\rho/2\le E_c<\rho(\Re S-Q/2),
 \qquad h\ge\Re S>N/2+Q/2>N/2.
\]
\end{proof}

\subsection{A finite Cayley--Hermite--Biehler inequality}

The identities \eqref{or:horizontal}--\eqref{or:radial-ratio} and the
next two source lemmas hold for every finite degree $n\ge2$,
$0<\rho<1$, $0<\alpha\le X\le1$, and $0<T\le1$ satisfying
\eqref{or:original-domain}--\eqref{or:critical-domain} at an actual
simple rightmost zero.  They also apply at the larger scales in the
next section.

\begin{lemma}[Finite polar source estimate]\label{or:HB}
For real $s$ suppose that
\begin{equation}\label{or:Cayley-entrance}
             \Re\frac{s}{1+ic_\rho s}>X.
\end{equation}
Define
\begin{equation}\label{or:source-definitions}
 \begin{gathered}
 W(s)=\sum_{F(u)=0}\frac{I(u)}{|s-(1+ic_\rho s)u|^2},\qquad
 C(s)=\sum_{F'(v)=0}\frac{I(v)}{|s-(1+ic_\rho s)v|^2},\\
 E(s)=\sum_{F(u)=0}\frac{I(u)}{|s-u|^2},\qquad
 \theta=\frac{n\rho}{2(1+c_\rho^2s^2)}.
 \end{gathered}
\end{equation}
Then
\begin{equation}\label{or:HB-inequality}
         W\ge E,\qquad \theta W\le(C+\rho)(W+\theta).
\end{equation}
In particular, if an upper bound $k$ for $C+\rho$ satisfies $k<\theta$,
then $W\le k/(1-k/\theta)$.
\end{lemma}
\begin{proof}
The maps
\[
 w=\frac u{1-ic_\rho u},\qquad
 u=\frac w{1+ic_\rho w},\qquad d_p(w)=1+ic_\rho w
\]
send the disk in \eqref{or:original-domain} to the closed lower
half-plane, because
\[
       -\Im w=\frac{I(u)}{|1-ic_\rho u|^2}.
\]
Consider the finite polynomials
\begin{equation}\label{or:polar-polynomials}
 G(w)=d_p(w)^nF(w/d_p(w)),\qquad
 J_p(w)=d_p(w)G'(w)-nic_\rho G(w)
       =d_p(w)^{n-1}F'(w/d_p(w)).
\end{equation}
A root $u=-i/c_\rho$ makes its
transformed linear factor constant.  It has $I(u)=0$ and contributes
zero to the sums in \eqref{or:source-definitions}.  All other labels
and multiplicities are preserved.  The coefficient $n$ in
\eqref{or:polar-polynomials} is the original degree, even if $G$ has
smaller degree.  Both $G$ and $J_p$ have their finite roots in the
closed lower half-plane.

By \eqref{or:Cayley-entrance}, $G(s)\ne0$.  Put
\[
 g=G'/G,\qquad a(w)=-nic_\rho/d_p(w),\qquad
 H_s(w)=G'(w)+a(s)G(w).
\]
Direct subtraction in the Cayley map proves
\[
 -\Im g(s)=W,\quad -\Im(J_p'/J_p)(s)=C,\quad
 -\Im a(s)=\theta,
 \qquad |g(s)+a(s)|\ge W+\theta>0.
\]
In particular, $H_s(s)$ and $J_p(s)$ are nonzero.  Differentiating
the polar identity before freezing $a$ gives
\begin{equation}\label{or:frozen-polar}
 \frac{H_s'}{H_s}(s)
  =\frac{J_p'}{J_p}(s)-\frac{ic_\rho}{d_p(s)}
                      -\frac{a'(s)}{g(s)+a(s)}.
\end{equation}
Here $|a'(s)|=c_\rho\theta$ and $|ic_\rho/d_p(s)|\le c_\rho$.
The last two terms therefore each have modulus at most $c_\rho$,
and
\begin{equation}\label{or:frozen-upper}
                   -\Im(H_s'/H_s)(s)\le C+\rho.
\end{equation}

The real-rootedness needed below follows from a finite argument.
If a degree $m$ polynomial has all roots strictly below
the real axis, a continuous argument of its value on the real axis
decreases strictly through a total angle $m\pi$: its derivative is
the negative sum of the positive Poisson kernels of the roots.
After any fixed rotation, each crossing of an integer multiple of
$\pi$ is a real zero of its imaginary part.  These crossings account
for the degree of that imaginary part, which is $m$, or $m-1$ when
its leading coefficient vanishes.  Hence that imaginary part is
real-rooted.  Moving roots a small distance downward and taking a
coefficient limit proves the assertion for roots on the boundary as
well, allowing repeated roots and degree drop.  A nonzero coefficient
limit of real-rooted polynomials is real-rooted, with roots allowed to
escape to infinity.

Rotate $G$ so that $H_s(s)>0$ is real, and write $G=U+iV$.  Let
$p=\Re a(s)$ and
\[
                       B=V'+pV-\theta U=\Im H_s.
\]
At the selected point $s$,
\[
 B(s)=0,\qquad H_s(s)V(s)=(W+\theta)|G(s)|^2>0,\qquad
 U'V-UV'=W|G|^2.
\]
In particular $V$ is not identically zero.  By the preceding argument
it is real-rooted, so
\[
 V'^2-VV''\ge0.
\]
This last inequality follows directly by logarithmic differentiation
away from the zeros, and by continuity at them.  The exact Wronskian
identity
\[
 BV'-B'V=\theta(U'V-UV')+V'^2-VV''
\]
now gives
\[
 \theta W\le[-\Im(H_s'/H_s)(s)](W+\theta)
             \le(C+\rho)(W+\theta).
\]
Finally, for $u=r-it$,
\begin{equation}\label{or:denominator-expansion}
 |s-(1+ic_\rho s)u|^2
   =|s-u|^2-\rho s^2t+\frac{\rho^2s^2}{4}|u|^2
   \le |s-u|^2,
\end{equation}
since $t\ge\rho|u|^2/2$.  Termwise comparison gives $W\ge E$.
Rearranging with $k<\theta$ gives the solved form.
\end{proof}

We will also use the unsolved form.  If $W\ge b_0>0$, then
\begin{equation}\label{or:parallel-source}
                   C+\rho\ge\frac{\theta b_0}{\theta+b_0}.
\end{equation}
The function $uv/(u+v)$ is increasing in each positive argument.  Thus
\eqref{or:parallel-source} remains useful without an upper bound on
$h$ or a positive solved denominator.

\subsection{The critical-label comparison}

For $s\ge\alpha$, the critical domain yields the exact identity
\begin{align}
 &|s-(1+ic_\rho s)v|^2
       -(1-\rho^2s^2/2)|s-v|^2\notag\\
 &\quad=\rho s^2\left(\frac\rho2|v-\alpha|^2-t\right)
 +\frac{\rho^2s^2}{2}(s-\alpha)(s+\alpha-2r)
 +\frac{\rho^2s^2}{4}|v|^2>0.\label{or:quadratic-polar}
\end{align}
All three terms have their stated sign on the whole unbounded critical
domain.  In particular, if $\rho\le r_+$ and $s\le s_+$, with
$r_+^2s_+^2/2<1$, set
\begin{equation}\label{or:polar-constants}
 f=\frac1{1-r_+^2s_+^2/2},\qquad
 D_p=1+(r_+s_+/2)^2.
\end{equation}
Then
\[
 C<f C_{\rm real},\qquad
 C_{\rm real}=\sum_v\frac{I(v)}{|s-v|^2},\qquad
 \theta\ge\frac{n\rho}{2D_p}.
\]

\begin{lemma}[A negative square-moment term]\label{or:critical-comparison}
For $s=X+\delta$ with $\delta\ge0$, put
\begin{equation}\label{or:kappa}
 \kappa_X(T,\delta)=\frac{X^2}{2}
 -\frac{X\max(0,T^2-X\delta-\delta^2)^2}
             {4(T^2+\delta^2)(X/2+\delta)}.
\end{equation}
Then
\begin{equation}\label{or:critical-comparison-formula}
 C_{\rm real}<\rho\{X\Re S-\kappa_X(T,\delta)Q\}.
\end{equation}
For $X=1$, the coefficient decreases with $T$, and is nonnegative for
$0<T\le1$, $\delta\ge0$.  For $T=1$ and
$\delta\in\{3/10,1/2,71/100\}$ there is also the bound
\begin{equation}\label{or:small-X-comparison}
 C_{\rm real}<m\rho\Re S\quad(0<X\le m),\qquad
                         m=(1-\delta)^2.
\end{equation}
\end{lemma}
\begin{proof}
Let $A=X-r>X/2$, $B=A+\delta$, $b_0=X/2+\delta$, and
$L=T^2+\delta^2$.  Use \eqref{or:inset-critical} and clear the
positive denominators of one summand in
\eqref{or:critical-comparison-formula}.  The coefficient available
after division by $X$ is
\[
 \frac X2+(A-X/2)
       \frac{2\delta A+\delta^2+2Tt-T^2}{(A+\delta)^2+t^2}.
\]
If the second numerator is nonnegative there is no loss.  Otherwise,
with $z_0=2\delta A+\delta^2-T^2<0$,
\[
 \frac{z_0+2Tt}{B^2+t^2}-\frac{z_0}{B^2}
       =\frac{2TtB^2-z_0t^2}{B^2(B^2+t^2)}\ge0.
\]
It suffices to take $t=0$.  The possible loss is
\[
 \frac{(B-b_0)(L-2\delta B)}{B^2}.
\]
As a function of $w=1/B\in[0,1/b_0]$ this is
$(1-b_0w)(Lw-2\delta)$.  Its maximum is
$\max(0,L-2\delta b_0)^2/(4Lb_0)$, by completing the square; when
$L-2\delta b_0\le0$ its maximum is zero at $w=1/b_0$.
This proves the stated coefficient.  For $X=1$ the subtracted term
increases with $T^2$ on its active interval.  Its numerator is at
most $T^4\le T^2\le(1+2\delta)(T^2+\delta^2)$, which also proves
nonnegativity of $\kappa_1$.

For \eqref{or:small-X-comparison}, put $T=1$.  The critical inset
divided by $\rho$ is less than $XA-X^2/2$.  If $A\ge m$ and $X\le m$,
this is at most $mA-m^2/2$.  The required comparison then reduces to
a quadratic in $A$ whose discriminant is
$[(1-\delta)^2-m][(1+\delta)^2+m]=0$; its remaining $t$ and $t^2$
coefficients are nonnegative.  If $A\le m$, use
$XA-X^2/2\le A^2/2$.  The constant term after clearing denominators is
\[
 P(A)=2m(A+\delta)^2-A(A^2+1).
\]
Its derivative is $-3A^2+4mA+4m\delta-1$.  In each of the three
specified cases,
$16m^2+12(4m\delta-1)<0$, so $P$ decreases on the real line.
Furthermore $P(m)=m\delta^4>0$.  The other terms are nonnegative
($2At$ and $(2m-A)t^2$ after the same division), proving the assertion.
\end{proof}

For a closed parameter cell suppose that
$C_{\rm real}<\rho h A$, where $0<A\le\overline A$, and that
$h\ge H$, $h/n\le U$.  With the constants in
\eqref{or:polar-constants}, if
\begin{equation}\label{or:HB-cell-denominator}
              B=2D_p(f\overline A U+1/n_0)<1,
\end{equation}
then Lemma~\ref{or:HB} gives
\begin{equation}\label{or:HB-cell}
                  \frac Wh<\frac{r_+(f\overline A+1/H)}{1-B}.
\end{equation}
If a ceiling $h<H_+$ has previously been proved, we may instead use
$B=2D_p(f\overline A H_++1)/n_0$.  A source lower bound $W\ge b_0$
gives, without either ceiling,
\begin{equation}\label{or:direct-K}
 \rho h A>K(\rho),\qquad
 K(\rho)=\frac1f\left\{\frac{n_0\rho b_0}
                         {n_0\rho+2D_p b_0}-\rho\right\}.
\end{equation}
The expression in braces is concave in $\rho$, since its second
derivative is $-4n_0^2D_pb_0^2/(n_0\rho+2D_pb_0)^3<0$.
It increases with $n_0$ when the original degree is retained.  Thus
the smaller endpoint value on a closed scale interval is a valid
lower bound for the entire interval.  In each application $K>0$,
so $A>0$ and $h>K/(\rho A)$.

\subsection{Moving sources and a uniform crescent gap}

When $X=1$ and $I=I(\beta)>0$, set
\begin{equation}\label{or:moving-test}
 D_0=1-c_\rho(T+I),\qquad s=\frac{1+I}{D_0}.
\end{equation}
The contribution of the actual zero $\beta$ is exactly
\begin{equation}\label{or:moving-weight}
 \frac{I}{|s-(1+ic_\rho s)\beta|^2}
       =\frac{D_0}{2I}\ge\frac{(1-\rho)^2}{2T}.
\end{equation}
Indeed $D_0=|1-ic_\rho\beta|^2$, and expansion gives
$D_0|s-(1+ic_\rho s)\beta|^2=2I^2$.
Moreover $(1-\rho)^2\le D_0<1$, $1<s<2/(1-\rho)^2<3$, and
\begin{align}
 D_0^2(s-1-c_\rho^2s^2)
 &=I\{1+2c_\rho(1-T)+c_\rho^2(T^2-4T-1)
                      +2c_\rho^3(1+T^2)\}>0.\label{or:moving-domain}
\end{align}
The bracket is at least $1-4c_\rho^2$: after subtraction the
remaining terms are
$2c_\rho(1-T)+c_\rho^2(T-1)(T-3)+2c_\rho^3(1+T^2)$.
Thus the test satisfies \eqref{or:Cayley-entrance}.

On a cell $\rho\le r_+$, $T\le t_+$ put
\begin{equation}\label{or:moving-field}
 b_m=\frac{(1-r_+)^2}{2t_+},\quad
 f_m=(1-9r_+^2/2)^{-1},\quad D_m=1+9r_+^2/4,\quad
 H_m=\frac{n_0b_m/(n_0r_++2D_mb_m)-1}{f_m}.
\end{equation}
Since $C<f_m\rho h$, the unsolved inequality
\eqref{or:parallel-source} proves $h>H_m$.  At zero inset it remains
valid with $h\ge H_m$.
At $X=1$ and $I(\beta)=0$,
\[
             T=\frac\rho2(1+T^2)\le\rho<1.
\]
To justify that endpoint, multiply every physical zero and the physical
witness by $1-\epsilon$.  For a sufficiently small positive $\epsilon$
all the finitely many strict holes persist, and $T'<1$ in the
$X=1$ normalization, where
\begin{equation}\label{or:homothety}
 \rho'=(1-\epsilon)\rho,\quad
 T'=T+\frac{\epsilon}{(1-\epsilon)\rho},\quad
 \beta'-v'=\beta-v,\quad S'=S,\quad h'=h.
\end{equation}
The scaled source has positive inset.  Applying the formula with a
depth endpoint tending down to $t_+$ and taking the rational limit
gives $h\ge H_m$.

We also need a distance bound which keeps the actual depth.  The lower
intersection of the hole circle with $r=\alpha/2$ has depth
\[
             \eta=\frac{1-\sqrt{1-\rho^2\alpha^2/4}}{\rho}.
\]
The inequality
\begin{equation}\label{or:height-gate}
 \frac2{125}-\rho_*\left(\frac14+\frac1{15625}\right)>0
\end{equation}
gives $0\le\eta<\alpha^2/125$ by squaring two positive sides.
To bound the distance on the critical crescent, write
$u=r-\alpha\le-\alpha/2$ on the lower hole arc.  The scalar
product with the vector $(X-\alpha,T-1/\rho)$ is
\[
 (X-\alpha)u+(1/\rho-T)\sqrt{\rho^{-2}-u^2}.
\]
Its derivative is positive for $u\le0$.  Hence the nearest permitted
arc point has the largest real coordinate.  The upper arc and exterior
radial segments cannot be nearer; if $T<\eta$, the horizontal
separation alone suffices.  We obtain
\begin{align}
 X=1&:\quad d^2>\frac14+\max(T-1/125,0)^2,\label{or:gap-X}\\
 T=1&:\quad d^2>G(X):=1+\frac{117}{500}X^2.\label{or:gap-T}
\end{align}
For the second line, the stronger lower bound is
$(X-\alpha/2)^2+(1-\alpha^2/125)^2
 \ge X^2/4+(1-X^2/125)^2
 =G(X)+X^4/15625$.

\subsection{From inverse contours to source moments}

We use the inverse and symmetric contour estimates of
Lemmas~\ref{lt:inverse-log} and~\ref{lt:quarter-pi},
and prove two refinements for one-sided strips.

Suppose $Q/h\le Y$, $1/d\le m$, and $h\ge H$.  For $0<mR<1$ put
\begin{equation}\label{or:inverse-parameters}
 \begin{gathered}
 \lambda=\frac{YR}{1-mR},\qquad
 e_R=Y\left\{\sum_{j=2}^{16}\frac{m^{j-2}R^j}{j}
                  +\frac{m^{15}R^{17}}{17(1-mR)}\right\},\\
 \widehat R=R-e_R,\qquad
 B_R=\frac{Y}{(1-mR)^2(1-\lambda)^3}.
 \end{gathered}
\end{equation}
The simpler upper bound
$e_R=Y\{R^2/2+mR^3/[3(1-mR)]\}$ is also permissible and will be
used for the base range.  These bound the complete logarithmic series
of the actual critical product; the last displayed term includes
every omitted order.  For example
\[
 \frac1S\log\frac{F'(\beta+z)}{F'(\beta)}=z+E_R(z),\qquad
 |E_R(z)|\le e_R\quad(|z|\le R),
\]
and the derivative differs from $1$ by at most $\lambda$.
If $\lambda<1$ and $\widehat R>0$, this logarithm is injective and
its image contains the disk of radius $\widehat R$.
In the unscaled coordinate
$w=\log(F'(\beta+z)/F'(\beta))$, the inverse position
$\phi(w)=\beta+z$ is defined on $|w|<h\widehat R$, with
$\phi(0)=\beta$ and $F'(\phi(w))/F'(\beta)=e^w$.
Set $A(w)=S\phi'(w)$.  The normalized primitive
$SF(\phi(w))/F'(\beta)$ has derivative $e^w A(w)$, and
\[
        \Re A\ge(1+\lambda)^{-1},\quad
        |A|\le(1-\lambda)^{-1},\quad |A'|\le B_R/h.
\]

Choose $0<\nu<1$, $W_0>0$, and rational bounds
\begin{equation}\label{or:contour-gates}
 \begin{gathered}
 B_R/H<\epsilon_0,\quad \nu B_R\widehat R<g_0,\quad
 0<p_0\le(1+\lambda)^{-1}<1<(1-\lambda)^{-1}\le a_0,\\
 e_0>e^{-W_0},\qquad
 W_0^2<(1-\nu^2)(H\widehat R)^2,\\
 \epsilon_0(1+e_0)+g_0e_0<1-a_0e_0,\qquad
 \epsilon_0(1+e_0)+g_0e_0^2<p_0-e_0.
 \end{gathered}
\end{equation}
For an odd-$\pi$ contour require
\begin{equation}\label{or:odd-gate}
      \epsilon_0+g_0e_0<p_0,\qquad \nu H\widehat R>\pi.
\end{equation}
For a quarter-period contour replace these requirements by
\begin{equation}\label{or:quarter-gate}
 \epsilon_0+g_0e_0<\frac{p_0}{1+p_0+a_0},\qquad
                        \nu H\widehat R>\pi/2.
\end{equation}
These are conditions on the containing disk, all corners, both vertical
sides, both horizontal sides, and the amplitude homotopy.  The resulting
number $L_R$ of actual original roots in $|u-\beta|<R$ obeys
\begin{equation}\label{or:counts}
 \frac{L_R}{h}>\frac{\nu\widehat R}{\pi}-\frac{\ell_0}{h},
 \qquad
 \ell_0=2\ \hbox{in \eqref{or:odd-gate}},\quad
 \ell_0=3/2\ \hbox{in \eqref{or:quarter-gate}}.
\end{equation}
The trimmed odd height is the largest odd multiple of $\pi$ not
exceeding the available height, including equality.  The full side
estimates already cover such an endpoint.  A one-sided odd-$\pi$ strip
between the fractions $\gamma$ and $\nu$ of the available height,
$0<\gamma<\nu$, gives
\begin{equation}\label{or:one-sided-count}
                    \frac{L_{\gamma,\nu}}h>
           \frac{(\nu-\gamma)\widehat R}{2\pi}-\frac2h.
\end{equation}

\begin{lemma}[One-sided inverse strips]\label{or:one-sided}
Under \eqref{or:contour-gates} and \eqref{or:odd-gate}, the upper
strip $\gamma h\widehat R<\Im w<\nu h\widehat R$ and its reflected
lower strip each have the count \eqref{or:one-sided-count}.
They exclude the source $\beta$.  Writing
$S/h=c_0+is_0$ with $c_0>0$, each upper-strip root lies above
$\beta$ by more than
\begin{equation}\label{or:strip-displacement}
                c_0\gamma\widehat R-e_R-|s_0|W_0/h;
\end{equation}
each lower-strip root lies below it by more than the same quantity.
\end{lemma}
\begin{proof}
Put $\mathcal J(w)=SF(\phi(w))/F'(\beta)$.
The estimates preceding \eqref{or:contour-gates} give
\[
 A(0)=1,\quad \Re A\ge p_0,\quad |A|\le a_0,\quad
 |A'|<\epsilon_0,\quad
 \mathcal J(w)=\int_0^w e^zA(z)\,dz.
\]
All corners of the containing rectangle lie in the inverse disk.
Integration by parts along the entire path
$0\longrightarrow-W_0\longrightarrow-W_0+i\eta
 \longrightarrow x+i\eta$ gives
\[
 \mathcal J(w)=e^wA(w)-1-\mathcal E(w),\qquad
 |\mathcal E(w)|<\epsilon_0(1+e^x)+g_0e^{-W_0}.
\]
The vertical segment contributes at most
$(B_R/h)|\eta|e^{-W_0}\le\nu B_R\widehat R e^{-W_0}$;
this gives the $g_0e^{-W_0}$ term.
On the left side the modulus of $e^wA(w)-1$ is at least
$1-a_0e^{-W_0}$, and on the right side it is at least
$p_0e^{W_0}-1$.  The two inequalities in the last line of
\eqref{or:contour-gates} control these whole sides, respectively.
On every horizontal line $\eta=(2k+1)\pi$ the modulus is at least
$1+p_0e^x$.  It dominates the error, because
\[
 \epsilon_0(1+e^x)+g_0e^{-W_0}
 \le(\epsilon_0+g_0e^{-W_0})(1+e^x)
 <p_0(1+e^x)\le1+p_0e^x.
\]
Thus Rouch\'e applies on every subrectangle bounded by two such
horizontal lines.  The homotopy
$A_t=1+t(A-1)$ preserves $\Re A_t\ge p_0$ and $|A_t|\le a_0$,
so $e^wA_t(w)-1$ stays nonzero on this boundary.
The zero count is therefore exactly that of $e^w-1$.

For the upper strip take the smallest odd multiple $L$ of $\pi$
at least $\gamma h\widehat R$, and the largest odd multiple $U$
at most $\nu h\widehat R$.  If $L<U$, the subrectangle contains
exactly $(U-L)/(2\pi)$ roots, and
\[
 \frac{U-L}{2\pi}>
            \frac{(\nu-\gamma)h\widehat R}{2\pi}-2.
\]
If the right side is negative, the claimed lower count is automatic;
if it is nonnegative, the strip width is at least $4\pi$ and
$L<U$ necessarily.  The same argument applies to $-U,-L$.
The boundaries contain no zeros, including when an endpoint is exactly
an odd multiple of $\pi$.  The inverse is injective and its image
is critical-free, so all counted roots are distinct simple actual
roots.  Since $\gamma>0$, neither strip contains $w=0$.
Finally
\[
             |\phi(w)-\beta-w/S|\le e_R,
 \qquad \Im(w/S)=\frac{c_0\Im w-s_0\Re w}{h},
\]
which proves \eqref{or:strip-displacement} and its lower-strip
counterpart.
\end{proof}

\begin{lemma}[Real widths and an upper root]\label{or:narrow-strip}
If $|\Im S|<J$, the roots counted in \eqref{or:counts} satisfy
\begin{equation}\label{or:width}
 |\Re(u-\beta)|\le\ell_R:=\min\left(R,
         \frac{W_0}{H}+\frac{J\nu\widehat R}{H}+e_R\right).
\end{equation}
For a row satisfying \eqref{or:odd-gate}, a positive upper strip
also supplies an actual root above $\beta$
by more than
\begin{equation}\label{or:depth}
 \left(1-\frac{J^2}{H^2}\right)\nu\widehat R-e_R
       -\frac{4\pi}{H}-\frac{JW_0}{H^2},
\end{equation}
whenever this expression is positive.  If it exceeds the depth $T$,
the original disk condition is contradicted.
\end{lemma}
\begin{proof}
In the counting rectangle $|\Re w|\le W_0$,
$|\Im w|\le\nu h\widehat R$.  Since
\[
 \Re(w/S)=\frac{c_0\Re w+s_0\Im w}{h},\qquad
 |s_0|<J/h,
\]
the inverse error bound gives \eqref{or:width}.
For the upper root let $U$ be the largest odd multiple of $\pi$
not exceeding $\nu h\widehat R$.
The adjacent rectangle $U-2\pi<\Im w<U$ contains exactly one
actual root by the complete boundary argument of
Lemma~\ref{or:one-sided}.  When \eqref{or:depth} is positive,
$\nu H\widehat R>4\pi$ and $J/H<1$, so that this rectangle has
positive height above zero.  Its root has
$\Im w>\nu h\widehat R-4\pi$.
Using $c_0=\sqrt{1-s_0^2}\ge1-(J/H)^2$, $c_0\le1$, and $h\ge H$
in the imaginary inverse displacement gives exactly
\eqref{or:depth}.  Every original zero has nonpositive imaginary
part in the normalized coordinates, which proves the last assertion.
\end{proof}

We now convert these counts into lower bounds for $E(s)$.  Suppose
$T\in[t_-,t_+]$, $\rho\le r_+$, and $s=X+\delta$.  A counted region
with radius $R$ and real width $\ell$ has $|\Re u|\le1$ in every
application below.  For $X=1$ this follows from rightmostness and
$\ell<2$; for $T=1$ it follows from $R<1$ and $0<X\le1$.
Its whole-region lower weight is
\begin{equation}\label{or:weight}
 \omega(R,\ell;\delta,t_-,t_+,r_+)
 =\max\left(0,\min_{t\in\{\max(0,t_--R),\,t_++R\}}
 \frac{t-r_+(1+t^2)/2}{(\delta+\ell)^2+t^2}\right).
\end{equation}
The derivative numerator of the quotient, writing
$A=(\delta+\ell)^2>0$, is
$A-t^2+r_+(1-A)t$.  Its two roots have opposite signs, and its
positive root is a maximum.  Thus the minimum on a nonnegative
interval is at an endpoint.  If that minimum is negative we instead
use $I(u)\ge0$.

Take increasing radii $R_j$; replace their widths by running maxima
so the regions are nested.  Let $s_j$ be any increasing rational lower
bounds for $\nu_j\widehat R_j/\pi$, and replace the weights by the
maximum of all later weights.  This gives valid
$w_1\ge\cdots\ge w_k\ge0$ on the nested regions.  The source root
has a separate lower weight $b_0$.  Subtracting that root once from
each odd-$\pi$ count, and applying finite Abel summation, gives
\begin{equation}\label{or:Abel-source}
 \frac Eh>
     \sum_{j=1}^k w_j(s_j-s_{j-1})+\frac{b_0-3w_1}{h},
                         \qquad s_0=0.
\end{equation}
Indeed, if $M_j$ are the actual cumulative counts excluding $\beta$,
then $M_j>hs_j-3$ and
$\sum w\ge w_kM_k+\sum_{j<k}(w_j-w_{j+1})M_j$.
The coefficients of the cumulative counts are nonnegative, so the bound
remains valid when an initial $hs_j-3$ is negative.
If $b_0-3w_1<0$, replace its denominator
in \eqref{or:Abel-source} by $H$.  If it is nonnegative and $h<600$,
replace it by $600$; if no upper field is known, discard the
nonnegative term.  This accounts for both signs of the source
compensation.
For a one-sided strip the corresponding calculation is simpler:
if its roots have weight at least $a>0$ and $\beta$ has weight at
least $b_0$, then, since they are distinct from $\beta$,
\begin{equation}\label{or:strip-source-compensation}
 \frac Eh>
 a\frac{(\nu-\gamma)\widehat R}{2\pi}+\frac{b_0-2a}{h}.
\end{equation}
Thus $b_0\ge2a$ makes the last term nonnegative.

For quarter-period counts, one may instead use the cumulative lower
bounds $k_j=\nu_j\widehat R_j/\pi-3/(2H)$ directly, when they and
their successive increments are positive.  Then
\begin{equation}\label{or:Abel-quarter}
                    \frac Eh>\sum_jw_j(k_j-k_{j-1}),\qquad k_0=0.
\end{equation}

\subsection{The base range $0<\rho\le1/100$}

All constants in this subsection use $r_b=1/100$.  Set
\begin{equation}\label{or:base-polar}
 f_m=\frac{1001}{1000},\qquad f_X=f_T=\frac{10001}{10000},
                    D_b=1+9r_b^2/4.
\end{equation}
These dominate the quadratic polar factors for, respectively, the
moving tests $s<3$, the fixed test $s=101/100$, and the fixed test
$s=X+3/10\le13/10$.  The fixed inverse-Cayley conditions are
$(r_b(101/100)/2)^2<1/100$ and
$(r_b(13/10)/2)^2<3/10$.  Lemma~\ref{or:critical-comparison} gives
$C_{\rm real}<\rho h$ for the $X$ tests, and
$C_{\rm real}<\rho\max(X,49/100)h$ for the $T$ test.

There are $55$ inverse rows in this range.  The columns in
Table~\ref{or:base-contours} are
\[
 (Y,m,R,H,\nu,W_0,\epsilon_0,g_0,p_0,a_0,e_0).
\]
The simpler $e_R$ following \eqref{or:inverse-parameters} is used.
\begin{table}[ht]
\centering\scriptsize
\caption{All inverse rows for the base range.  Each $j$-row occurs for
every integer $6\le j\le30$.}\label{or:base-contours}
\begin{tabular}{@{}lrrrrrrrrrrr@{}}
\toprule
row&$Y$&$m$&$R$&$H$&$\nu$&$W_0$&$\epsilon_0$&$g_0$&$p_0$&$a_0$&$e_0$\\
\midrule
$X_1$&2&2&$3/20$&600&$99/100$&4&$1/25$&3&$7/10$&$7/4$&$1/50$\\
$X_D$&2&2&$13/100$&600&$99/100$&4&$1/25$&3&$7/10$&$7/4$&$1/50$\\
$X_{450,j}$&2&2&$j/200$&450&$23/25$&$W_j$&$11/100$&$g_j$&$7/10$&$7/4$&$e_j$\\
$X_{200,j}$&2&2&$j/200$&200&$9/10$&$W_j$&$11/100$&$g_j$&$7/10$&$7/4$&$e_j$\\
$X_{78}$&2&2&$3/20$&78&$97/100$&$11/5$&$3/10$&3&$7/10$&$7/4$&$1/9$\\
$T_{150}$&$25/11$&$100/99$&$3/20$&150&$9/10$&4&$1/6$&3&$7/10$&$7/4$&$1/50$\\
$T_{103}$&$25/11$&$100/99$&$3/20$&103&$9/10$&4&$1/6$&3&$7/10$&$7/4$&$1/50$\\
\bottomrule
\end{tabular}
\end{table}
Here $(W_j,g_j,e_j)=(2,1/5,1/7)$ for $6\le j\le9$, and
$(4,3,1/50)$ for $10\le j\le30$.  Only the two $j$-families use
quarter-period contours; all other rows use odd-$\pi$ contours.
Every row satisfies \eqref{or:contour-gates} and its appropriate
horizontal and height condition.  These are rational inequalities
apart from $e^{-W_0}$; the latter is bounded by the reciprocal of the
positive Taylor sum through degree $48$.  Thus Table~\ref{or:base-contours}
and the displayed formulae specify the finite verification.

\subsubsection{The base case $X=1$: the first ceiling}

Suppose $h\ge600$.  Then $d>1/2$ and $y_0<2$, so
$q<2/600<7/625$ and $\mu<2/600<1/200$.
Proposition~\ref{lt:fine200} therefore gives
\begin{equation}\label{or:base-high-phase}
 |\Im S|<4y_0<8,\qquad
 |\Im(S/h)|<1/70,\qquad \Re(S/h)>999/1000.
\end{equation}
The upper one-sided strip of row $X_1$, with $\gamma=4/5$, has
positive count, since
\[
 \frac{(99/100-4/5)\widehat R_1}{2(22/7)}-\frac2{600}>0.
\]
Its actual root lies above $\beta$ by more than
$(4/5)(999/1000)\widehat R_1-e_{R_1}-4/(70\cdot600)>67/1000$.
Thus $T>67/1000$.

First let $67/1000\le T\le3/10$.  Equation~\eqref{or:radial-ratio},
with $J=9$, gives $u=h/n<311/1000=:U$ for
every $N\ge6250$.  At $s=101/100$,
\[
       \kappa_1(T,1/100)\ge\kappa=\frac{29591}{63600}>0.
\]
As $y_0>u$, put $a(u)=f_X(1-\kappa u)$ and
$B(u)=2D_b(a(u)u+1/n_0)$.  On $0\le u\le U$ the factors
$1-\kappa u$, $1-2\kappa u$, and $1-B(u)$ are positive.
The derivative numerator of
$(a(u)+1/600)/(1-B(u))$ is
\[
 -f_X\kappa(1-2D_b/n_0)
 +2D_b\{f_X^2(1-\kappa u)^2+(f_X/600)(1-2\kappa u)\}.
\]
It has a positive lower bound obtained by replacing the first term
by $-f_X\kappa$ and $u$ by $U$ in the remaining terms.
Therefore the solved source estimate is valid throughout this
interval and
\begin{equation}\label{or:base-low-depth-upper}
                  W/h<r_b\frac{a(U)+1/600}{1-B(U)}<1/50.
\end{equation}

Use the lower one-sided strip of $X_D$, with $\gamma=7/25$, and set
\[
 \begin{split}
 \Delta&=\frac7{25}\frac{999}{1000}\widehat R_D-e_{R_D}
          -\frac4{70\cdot600}-\frac{3481}{5000}r_b>0,\\
 a_D(T)&=\frac{T+\Delta}{(T+157/1200)^2},\qquad
 b_D(T)=\frac{T-r_b(1+T^2)/2}{T^2+1/10000}.
 \end{split}
\]
These are full strip and source weights.  The strip depth estimate
comes from \eqref{or:base-high-phase}; its roots have
$|u|<\sqrt{1.09}+.13<59/50$ and
$|s-u|<\sqrt{T^2+.0001}+.13\le T+157/1200$.
The polynomial
\[
 [T-r_b(1+T^2)/2](T+157/1200)^2
             -2(T+\Delta)(T^2+1/10000)
\]
has all its Bernstein coefficients positive on
$[67/1000,3/10]$.  Hence $b_D(T)>2a_D(T)$ there.
Adding $\beta$ to the strip count in \eqref{or:one-sided-count}
compensates the loss $2/h$ and yields
\[
 \frac Eh>\frac{(99/100-7/25)\widehat R_D}{2(22/7)}a_D(T)>1/50.
\]
For the final comparison it suffices to check the two endpoints:
the only positive stationary point of $a_D$ is a maximum.
This contradicts \eqref{or:base-low-depth-upper}.

If $3/10\le T\le1$, equation~\eqref{or:radial-ratio} gives
$h/n<147/1000$.  Its source denominator is less than $3/10$, and
\[
 \frac Wh<
 \frac{r_b(f_X+1/600)}{1-2D_b(f_X(147/1000)+1/n_0)}<2/125.
\]
The symmetric $X_1$ count has lower weight
\[
       \frac{T-3/20-5r_b/4}{(T+901/6000)^2}.
\]
Indeed $|u|<\sqrt2+3/20<\sqrt{5/2}$ and the displayed denominator
bounds the complete disk.  Multiplication by
$(99/100)(7/22)\widehat R_1-2/600$ gives a lower bound exceeding
$1/50$ at both endpoints, and therefore on the whole interval, again
by its single-maximum derivative.  Thus $h\ge600$ is impossible.

\subsubsection{The base case $X=1$: the remaining fields}

We now have $h<600$.  Lemma~\ref{or:zero-inset} excludes $I(\beta)=0$,
since $N/2\ge3125$.
Apply the moving-source estimate at this positive inset.
The following table gives entire-domain denominator and depth bounds:
\begin{equation}\label{or:base-moving-table}
\begin{array}{c|c|c}
 \text{known ceiling }h<H_+& (C+\rho)/\theta<B_+&T>t_+\\ \hline
 600&193/1000&13/200\\
 450&29/200&9/100\\
 200&13/200&21/100
\end{array}
\end{equation}
For each row the two exact checks are
\[
 \frac{2D_b(H_+f_m+1)}{6251}<B_+,\qquad
 \frac{(1-r_b)^2}{2t_+}>
                   \frac{r_b(H_+f_m+1)}{1-B_+}.
\]
Thus they hold for every positive $\rho\le r_b$.

Suppose first $h\ge450$ and use all $25$ rows $X_{450,j}$.
For each radius, apply \eqref{or:weight} with $\ell=R$ and
$\delta=1/100$ on each of the two closed depth intervals
$[13/200,1/5]$ and $[1/5,1]$.  Use one of these depth intervals
for the entire nested family.  The cumulative bounds
\[
 k_j=\nu_j(7/22)\widehat R_j-\frac{3}{2H}
\]
and their successive increments are positive, and the resulting
weights decrease.  Formula~\eqref{or:Abel-quarter} applies.
After excluding this field, repeat with $H=200$ and depth intervals
$[9/100,1/5]$, $[1/5,1]$.  The exact inequalities, rounded outwards
here to shorter rational numbers, are
\begin{equation}\label{or:base-middle-table}
\begin{array}{c|c|c|c}
 H&[t_-,t_+]&E/h\text{ exceeds}&W/h\text{ is less than}\\ \hline
 450&[.065,.2]&.042&.0125\\
 450&[.2,1]&.029&.0125\\
 200&[.09,.2]&.041&.0118\\
 200&[.2,1]&.0245&.0118
\end{array}
\end{equation}
The source upper in each row is exactly
$r_b(f_X+1/H)/(1-B_+)$ with the previously established $B_+$.
These comparisons prove successively $h<450$, $T>.09$, $h<200$,
and $T>.21$.

On $[.21,1]$ the source weight is at least
\[
                   b_X=\frac{9900}{10001}.
\]
To check the whole interval, the numerator difference after clearing
the positive denominators factors as $(1-T)(AT-B)$, where
$A=1-(r_b/2)(1-1/10000)$,
$B=(r_b/2)(1-1/10000)+1/10000$, and $A(.21)>B>0$.
Let
$a_{78}=\omega(3/20,3/20;1/100,.21,1,r_b)$.
If $h\ge78$, the row $X_{78}$ and source compensation give
\[
 \frac Eh>
 a_{78}\frac{97}{100}\frac7{22}\widehat R_{78}
            +\frac{\min(0,b_X-3a_{78})}{78}>.01176,
 \qquad
 \frac Wh<\frac{r_b(f_X+1/78)}{1-.065}<.01084.
\]
Thus $h<78$.  Finally,
\[
              W<\frac{r_b(78f_X+1)}{1-.065}<.846<b_X\le E,
\]
which excludes the entire base case $X=1$.

\subsubsection{The base case $T=1$}

A direct physical comparison gives $d>99/100$.  In fact, with
$d_0=99/100$,
\[
 (1-d_0\rho)^2-|b-a|^2
 \ge\rho\{2(1-d_0)-(2-d_0^2)\rho\}>0,
\]
and $|a-\xi|>1$ then implies $|b-\xi|>d_0\rho$.
If $h\ge150$, this first gives $dh>100$, so
Proposition~\ref{mf:D100} yields $y_0<25/11$.  The bounds
$q<25/(11\cdot150)<9/400$ and
$\mu<1/(.99\cdot150)<1/100$ permit Proposition~\ref{lt:wideD100},
giving $|\Im S|<(23/5)y_0<11$.
Equation~\eqref{or:radial-ratio} implies $h/n<21/250$.
The row $T_{150}$ then gives
\[
 \frac Eh>.01599>.01211>
 \frac{r_b(f_T+1/150)}{1-2D_b(f_T(21/250)+1/n_0)}>\frac Wh.
\]
Thus $h<150$, whence $(C+\rho)/\theta<1/20$.

If $h\ge103$, the bound $103(.99)>100$ permits
Proposition~\ref{mf:D100}.
The row $T_{103}$ gives
\[
 \frac Eh>.01145>.01063>
              \frac{r_b(f_T+1/103)}{.95}>\frac Wh.
\]
Hence $h<103$.  The source $\beta$ itself always has weight at least
$b_T=99/109$.  For $X\le.77$, the uniform small-$X$ comparison
gives
\[
            W<\frac{r_b(103f_T(.77)+1)}{.95}<.846<b_T,
\]
a contradiction.

It remains to consider $X>.77$.  If $h\ge79$,
\eqref{or:horizontal} and $d>.99$ give
\begin{equation}\label{or:base-quartic-entry}
             q<\frac2{.77\cdot79}<\frac{33}{1000},
             \qquad \mu<\frac1{.99\cdot79}<\frac8{625}.
\end{equation}
The quartic construction of Proposition~\ref{lt:quartic} supplies two
distinct simple actual roots at distance less than $15/(2h)$ from
$\beta$.  For this construction, at $q_b=33/1000$, $\mu_b=8/625$,
$\nu_b=q_b\mu_b$,
use
\[
 (Z_0,M_0,E_0,P_0,Z_{\mathrm{out}},L_0,R_0,\eta_0)
   =(719/100,7/25,71/50,350,15/2,32/25,3/10,11/50).
\]
The center and path majorants of that construction are bounded by
\[
 \begin{split}
 Z(q)&=63/10+21q+133q^2+1030q^3+9000q^4,\\
 M(q)&=(63/10)q+45q^2+400q^3+3400q^4+86000q^5\\
     &\quad+320000q^6+1700000q^7+9100000q^8+39000000q^9.
 \end{split}
\]
They satisfy $Z(q_b)<Z_0$, $M(q_b)<M_0$,
$\mu_bZ_{\mathrm{out}}<1$, $Z_0+R_0<Z_{\mathrm{out}}$, and
\begin{align}
 \exp\left(M_0+\frac{\nu_bZ_0^3}{3(1-\mu_bZ_0)}\right)&<E_0,
 &\frac{Z_0^4E_0}{12(1-\mu_bZ_0)}&<P_0,\notag\\
 E_0(132200q_b^5+P_0\nu_b)&<(1-\eta_0)R_0,
 &1+q_bZ_{\mathrm{out}}+
 \frac{\nu_bZ_{\mathrm{out}}^2}{1-\mu_bZ_{\mathrm{out}}}&<L_0,
 \label{or:quartic-gates}\\
 \frac{e^{L_0R_0}-1-L_0R_0}{L_0}&<\eta_0R_0.&&\notag
\end{align}
The complete Gaussian residual is $132200q_b^5$, on its valid
domain $q_b<1/20$; the integrated actual cubic logarithmic tail
retains the factors $1/3$ and $1/4$ in
\eqref{or:quartic-gates}.  The last two residual bounds give strict
Rouch\'e inequalities on both entire circles.  Those circles avoid
the source and are disjoint because
\[
 6-21q_b-133q_b^2-1030q_b^3-9000q_b^4>R_0,
 \quad
 12(1-q_b-17q_b^2-137q_b^3-313q_b^4)>2R_0.
\]
Rouch\'e's theorem therefore gives the two distinct simple zeros of $F$.

Since $15/(2\cdot79)<19/200$, each companion has weight at least
\[
 \omega(19/200,19/200;3/10,1,1,r_b)
                         =\frac{8672039}{10840400}.
\]
Adding the source once gives
$E\ge b_T+2(8672039/10840400)>r_b(103f_T+1)/.95>W$.
Thus $h<79$, and the remaining source comparison
$W<r_b(79f_T+1)/.95<.843<b_T\le E$ is again impossible.
This proves Proposition~\ref{or:theorem} for all $0<\rho\le1/100$.

\subsection{The higher strip with $T=1$: entering the local estimates}

Assume throughout this subsection that
\begin{equation}\label{or:T-domain}
           1/100\le\rho\le\rho_*,\quad T=1,\quad 0<X\le1.
\end{equation}
The closed lower endpoint overlaps the base range.  By
\eqref{or:gap-T}, $d>\sqrt{G(X)}>1$.
For a test $s=X+\delta$ write
\[
 \kappa_\delta(X)=\kappa_X(1,\delta),\qquad
 A_\delta=X-\kappa_\delta(X)y_0,
 \quad b_\delta(X)=\frac{1-\rho_*(1+X^2)/2}{1+\delta^2}.
\]
On the active interval $1-X\delta-\delta^2\ge0$ the coefficient factors as
\begin{equation}\label{or:T-kappa-factor}
 \kappa_\delta(X)=
 \frac{X[X-(1-\delta)^2][X+(1+\delta)^2]}
                 {2(1+\delta^2)(X+2\delta)};
\end{equation}
outside it $\kappa_\delta(X)=X^2/2$.  With
$D_\delta(X)=1+\rho_*^2(X+\delta)^2/4$ and
$f_\delta(X)=(1-\rho_*^2(X+\delta)^2/2)^{-1}$, define
\begin{equation}\label{or:T-source-function}
 K_\delta(r,X)=\frac1{f_\delta(X)}
 \left\{\frac{n_0r b_\delta(X)}{n_0r+2D_\delta(X)b_\delta(X)}-r\right\}.
\end{equation}
Equation~\eqref{or:direct-K} gives
$\rho hA_\delta>\min_{r\in\{1/100,\rho_*\}}K_\delta(r,X)$.
All tests used here have $1/200\le\delta\le1$; the inverse-Cayley
condition follows from $\rho_*^2<1/200$.

For $\delta=1/2$, replacing the source and polar factors by their
$X=1$ bounds gives
\begin{equation}\label{or:Kh}
 K_h=\min_{r\in\{1/100,\rho_*\}}K_{1/2}(r,1)>.6843,
 \qquad \rho h A>K_h,
\end{equation}
where $A=1/4$ if $X\le1/4$, by
\eqref{or:small-X-comparison}, and
$A=X-\kappa_{1/2}(X)y_0$ otherwise.  Here
\[
              \kappa_{1/2}(X)=\frac{X(4X-1)(4X+9)}{40(X+1)}.
\]
If $X\le2/5$, \eqref{or:Kh} gives
$D>h>K_h/((2/5)\rho_*)>25$.
For $2/5\le X\le1$, the one-variable inequality
\begin{equation}\label{or:T-first-square}
            4(227/200)\kappa_{1/2}-(X+\kappa_{1/2})^2>0
\end{equation}
implies $A(1+y_0)<227/200$ by completion of the square in $y_0$.
Multiplication of \eqref{or:Kh} and
\eqref{or:radial-unphased} then gives
$h^2>K_hN/[2(227/200)]>25^2$.
Thus Proposition~\ref{mf:D25} is applicable in all cases, and
\begin{equation}\label{or:T-ratio3}
                  y_0<3/d<3/\sqrt{G(X)}\le3.
\end{equation}

The same moment ratio sharpens the joint bound to
\begin{equation}\label{or:T-joint-gap}
                  \frac{A(1+y_0)}{G(X)}\le\frac{21}{20}.
\end{equation}
Here are whole-interval certificates for this assertion.  For
$X\le1/4$ use $(1/4)(1+3)=1$.  Put
$Y(X)=3(1-11X^2/100)$.  Squaring positive sides gives
$Y(X)>3/\sqrt{G(X)}$ on $0<X\le.49$; after removal of $X^2$ the
required polynomial is strictly positive.  The ratio
$\kappa_{1/2}(X)/X$ increases, and seven times its value at $.41$
is less than one.  Hence on $[1/4,.41]$ the product
$(X-\kappa_{1/2}y_0)(1+y_0)$ increases up to its allowed endpoint,
and the certificate is
\[
 (21/20)G(X)-(X-\kappa_{1/2}Y(X))(1+Y(X))>0.
\]
On $[.41,1]$ completion of the square gives the certificate
$4(21/20)G(X)\kappa_{1/2}-(X+\kappa_{1/2})^2>0$.
Consequently
\begin{equation}\label{or:T-D45}
                 D^2>\frac{K_hN}{2(21/20)}>45^2.
\end{equation}

Write $S/h=c_0+is_0$, $c_0>0$, $c_0^2+s_0^2=1$.
The universal envelope in Proposition~\ref{mf:envelope} is
\begin{equation}\label{or:T-matrix}
 q<\frac{21}{5}\mu^2+\frac{11}{2}\mu c_0
                         +\frac{c_0^2}{4}-\frac{s_0^2}{2}.
\end{equation}
If $|s_0|\ge.7$, then $c_0<18/25$, and the right side is at most
\[
 \frac{21/5}{45^2}+\frac{(11/2)(18/25)}{45}
                       +\frac{1-.7^2}{4}-\frac{.7^2}{2}<0.
\]
Thus $|\Im S|<.7h$.

Now take $\delta=3/5$.  On $[0,1]$ the coefficient is always active.
The rational functions
$K_{3/5}(1/100,X)-K_{3/5}(\rho_*,X)$ and
$K_{3/5}(\rho_*,X)$ are positive.  Write
\[
 K_{.6}(X)=K_{3/5}(\rho_*,X),\quad
 A_{.6}=X-\kappa_{3/5}(X)y_0,\quad
 R_{.6}(y_0)=.7+.3\rho_*+(199/200)y_0.
\]
Equation~\eqref{or:radial} gives $N\rho/2<hR_{.6}(y_0)$.
Combining this with $\rho hA_{.6}>K_{.6}(X)$ gives
\[
 h^2>\frac{NK_{.6}(X)}{2A_{.6}R_{.6}(y_0)},\qquad
 D^2>\frac{NK_{.6}(X)G(X)}{2A_{.6}R_{.6}(y_0)}.
\]
The two rectangle inequalities are
\begin{align}
 N_0K_{.6}(X)G(X)-2\cdot51^2 A_{.6}R_{.6}(y_0)&>0,
       &&0\le X\le1, 0\le y_0\le3,\label{or:T-D51-poly}\\
 N_0K_{.6}(X)-2\cdot45^2 A_{.6}R_{.6}(y_0)&>0,
       &&0\le X\le1, 0\le y_0\le3.\label{or:T-h45-poly}
\end{align}
They prove $D>51$ and $h>45$.  An exact finite positivity procedure
for these two rational quadratics and all the preceding one-variable
certificates is given below.  If $|s_0|\ge2/3$, then $c_0<3/4$, and
\[
 \frac{21/5}{51^2}+\frac{(11/2)(3/4)}{51}
                                  +\frac5{36}-\frac29<0.
\]
Therefore
\begin{equation}\label{or:T-prephase}
                      h>45,\quad D>51,\quad |\Im S|<(2/3)h.
\end{equation}

\subsubsection{Entry to the larger companion rectangle}

Set $\delta=.71$ and $K(X)=K_\delta(\rho_*,X)$.
The same endpoint comparison proves that $K(X)$ is the smaller
scale endpoint; also $K(X)\ge K_*:=K(1)>0$.
Using \eqref{or:T-prephase} in the exact radial identity gives
\[
 \frac{N\rho}{2}<hR,\qquad
 R=\rho X+\frac23(1-\rho)
                   +\left(1-\frac\rho2(1+X^2)\right)y_0.
\]
Consequently the same source gives
\begin{equation}\label{or:T-q-square}
 q^2<\frac{2y_0^2[X-\kappa_\delta(X)y_0]R}{NK(X)}.
\end{equation}
For $X\le.49$, $R$ decreases with $\rho$.  At $\rho=.01$ put
$\beta_0=33/50+X/100$, $I_0=199/200-X^2/200$.
If $r(X)=\kappa_\delta(X)/X$, the derivative in $y_0$ of
$y_0^2(X-\kappa_\delta y_0)(\beta_0+I_0y_0)$, divided by $Xy_0$,
is
\[
       2\beta_0+3(I_0-r\beta_0)y_0-4rI_0y_0^2.
\]
The ratio $r$ increases: after positive denominators are removed
its derivative numerator is
$\delta^4+6\delta^2+4\delta X+X^2+1>0$.
Let $r_+$ be its value at $.49$, and take the endpoint extrema of
$\beta_0,I_0$ on $[0,.49]$.  The last display is bounded below by
\[
 2\beta_-+3(I_--r_+\beta_+)y_0-4r_+I_+y_0^2.
\]
This concave quadratic is positive at $0$ and $3$.  It follows that
the permitted $y_0$ may be replaced by $Y(X)$ from above.  The entire
interval is then certified by
\begin{equation}\label{or:T-q-low-poly}
 N_0K(X)-\frac{2}{(.057)^2}Y(X)^2
       [X-\kappa_\delta(X)Y(X)](\beta_0+I_0Y(X))>0
                   \quad(0\le X\le.49).
\end{equation}
For $X\ge.49$ the active interval has
$\kappa_\delta(X)\ge(12/25)X^2$; its inactive value $X^2/2$ is
stronger.  Completion of the square gives
$y_0[X-\kappa_\delta(X)y_0]\le25/48$.
Also \eqref{or:T-ratio3} gives $y_0<73/25$.
Thus the right side of \eqref{or:T-q-square} is at most
\[
 \frac{2(25/48)(73/25)
       \{2/3+\rho_*/3+(199/200)(73/25)\}}{N_0K_*}<(.057)^2.
\]
We have proved
\begin{equation}\label{or:T-Q057-entry}
                            q<57/1000,\qquad \mu<1/51.
\end{equation}
Proposition~\ref{lt:W8} applies in the form \eqref{lt:Q057}.
Its two distinct actual roots
$u_\pm$ have normalized positions $w_\pm=S(u_\pm-\beta)$ satisfying
\begin{equation}\label{or:T-Q057-pair}
 |w_\pm|<8917/1000,\quad
 \left|\frac{w_++w_-}{2}\right|<2109/1000,\quad
                         |\Im S|<(37/5)y_0.
\end{equation}
The mean estimate uses the even part of the two centers at the same
moment parameters, as specified in that proposition.

\subsection{The higher $T=1$ strip: two field ceilings}

If $h\ge105$, then $D>100$ by \eqref{or:gap-T}.
Propositions~\ref{mf:D100} and \ref{lt:wideD100} give
\begin{equation}\label{or:T-first-phase}
         y_0<9/4,\quad q<9/400,\quad \mu<1/100,\quad
                         |\Im S|<(23/5)y_0.
\end{equation}
For a closed cell $\rho\in[p,r]$, $y_0\in[l,u]$, take
$J=(23/5)u$ and a rational $H\ge105$ with either $H=105$ or
\begin{equation}\label{or:T-field-H}
                  H\le\frac{N_0p/2-J}{r+u}.
\end{equation}
Then $h\ge H$.  Put
\begin{equation}\label{or:T-radial-cell}
             B_r=\frac{r/2+J/N_0}{1-r/2}.
\end{equation}
A positive rational $U$ is a valid bound for $h/n$ if at least one
of $U\ge u$, $U^2\ge B_r$, or $Ul\ge B_r$ holds; the last is
used only for $l>0$.  These follow respectively from
\eqref{or:horizontal}, \eqref{or:radial-ratio}, and
\eqref{or:radial-sharp}.

Retain the same-$X$ distance information.  A cap $0<X_+\le1$ is
valid if $X_+=1$ or
\begin{equation}\label{or:T-Xcap}
                         l^2G(X_+)\ge(9/4)^2.
\end{equation}
For the chosen $\delta$, certify
\begin{equation}\label{or:T-Abar}
 \overline A-X+\kappa_\delta(X)v>0
       \quad(0\le X\le X_+,\quad v\in\{l,u\}).
\end{equation}
Split the $X$ interval at $(1-\delta^2)/\delta$ to use the correct
active or inactive formula.  Since the expression is affine in
$v$, this verifies the whole $(X,y_0)$ cell.  The constants in
\eqref{or:polar-constants} use $s_+=1+\delta$ and $r_+=r$.
After checking \eqref{or:HB-cell-denominator},
\eqref{or:HB-cell} is the source upper bound.

For each contour use \eqref{or:inverse-parameters} with $Y=u$, $m=1$,
and set
\begin{equation}\label{or:automatic-contour}
 \epsilon_0=\frac{1001B_R}{1000H},\quad
 g_0=\frac{1001\nu B_R\widehat R}{1000},\quad
 p_0=(1+\lambda)^{-1},\quad a_0=(1-\lambda)^{-1},\quad
 e_0=\left(\sum_{j=0}^{80}\frac{W_0^j}{j!}\right)^{-1}.
\end{equation}
Verify \eqref{or:contour-gates} and \eqref{or:odd-gate}, using
$\pi<22/7$.  Use \eqref{or:width}, rounded upwards, and the weights
\eqref{or:weight} with $t_-=t_+=1$.  A valid source weight is
$b_0=(1-r)/(1+\delta^2)$.  The source lower bound is
\begin{equation}\label{or:T-cloud-lower}
  \sum_jw_j(s_j-s_{j-1})+\frac{\min(0,b_0-3w_1)}H.
\end{equation}
The source upper is \eqref{or:HB-cell}.  Table~\ref{or:T-cover}
specifies the two full covers; the rational choices of cells and
contours are in the electronic appendix.  On every cell of the first
cover, \eqref{or:T-cloud-lower} strictly exceeds that upper.  The
coverage and arithmetic procedures are specified in
\S\ref{or:finite-verification}.  This proves $h<105$.

Suppose next that $90\le h<105$.  Since $D>90$,
Proposition~\ref{lt:D70} gives, with $z=Qd/h$,
\[
 z<8/3,\quad |\Im S|<(21/4)y_0,\quad
                   3z^2-(D-91/10)z+2D>0.
\]
At $D=90$ this quadratic is negative at both $z=8/3$ and $z=49/20$.
These values decrease with $D$, since $2-z<0$.  The prior
$z<8/3$ excludes the larger root branch, so
\begin{equation}\label{or:T-ratio245}
                           z<49/20,\qquad y_0<49/20.
\end{equation}
Use the nested-count procedure with $H\ge90$, $J=(21/4)u$ and
$49/20$ in place of $9/4$ in \eqref{or:T-Xcap}.
It excludes $26$ closed cells of the second cover.  The remaining
nine cells use the already proved $h<105$ and the following pair
estimate.

\begin{lemma}[Two actual sources]\label{or:pair-source}
Suppose $T=1$, $0<\rho\le r<1$, $R>0$, and $0<m<1$.
Suppose two distinct actual zeros other than $\beta$ have
displacements $v_\pm=u_\pm-\beta$ satisfying
$|v_\pm|<R$, $|(v_++v_-)/2|<m$.  Suppose also that their individual
insets belong to $[3/4,6/5]$.  For a test
$s=X+\delta$ with $\delta>0$, put
\[
 I_\Sigma=2(1-r)(1-m)-rR^2,\qquad
 D_\Sigma=2(1+\delta^2)+4\sqrt{1+\delta^2}\,m+2R^2.
\]
Assume $I_\Sigma>0$.  Then
\begin{equation}\label{or:pair-lower}
        E>\frac{1-r}{1+\delta^2}
                    +\frac{(I_\Sigma/2+49/50)^2}{D_\Sigma}.
\end{equation}
\end{lemma}
\begin{proof}
Rightmostness gives $\Re v_\pm\le0$.  Expanding the inset,
\[
 I(\beta+v)=I(\beta)-(1-\rho)\Im v
                             -\rho X\Re v-\frac\rho2|v|^2.
\]
Since $I(\beta)\ge1-\rho$ and $|v_++v_-|<2m$,
\[
 I(u_+)+I(u_-)
 >2(1-\rho)(1-m)-\rho R^2
 \ge2(1-r)(1-m)-rR^2=I_\Sigma>0.
\]
The second inequality uses $2(1-m)+R^2>0$.
The squared distances satisfy
\[
 \begin{split}
 |s-u_+|^2+|s-u_-|^2
 &=2(1+\delta^2)
   -2\Re\bigl((\delta+i)\overline{v_++v_-}\bigr)
   +|v_+|^2+|v_-|^2\\
 &<D_\Sigma.
 \end{split}
\]
On $[3/4,6/5]$ one has $\sqrt t\ge t/2+49/100$.
Indeed $t-(t/2+49/100)^2$ is concave and positive at both endpoints.
Cauchy--Schwarz gives
\[
 \frac{I(u_+)}{|s-u_+|^2}+\frac{I(u_-)}{|s-u_-|^2}
 \ge\frac{(\sqrt{I(u_+)}+\sqrt{I(u_-)})^2}
              {|s-u_+|^2+|s-u_-|^2}
 >\frac{(I_\Sigma/2+49/50)^2}{D_\Sigma}.
\]
Adding the contribution of $\beta$ proves \eqref{or:pair-lower}.
\end{proof}

By \eqref{or:T-ratio245}, the field $h\ge90$ gives
$q<(49/20)/90<.05$, $\mu<1/90<.02$.
Proposition~\ref{lt:q05mu02} therefore supplies
\[
                          R=17/180,\qquad m=19/900.
\]
Here $R>0$, $0<m<1$, and, for each of the nine cells,
\[
 I_\Sigma\ge2(1-\rho_*)(1-m)-\rho_*R^2
             =\frac{367561059}{200000000}>0.
\]
The individual insets lie in $[3/4,6/5]$, since
\begin{equation}\label{or:pair-inset-gates}
          (1-r)(1-R)-rR^2/2>3/4,\qquad 1+R<6/5.
\end{equation}
For each pair cell verify \eqref{or:T-Abar} and
$B_{105}=2D_p(105f\overline A+1)/n_0<1$.  Its source upper is
\begin{equation}\label{or:T-pair-upper}
                  W<\frac{r(105f\overline A+1)}{1-B_{105}}.
\end{equation}
The lower bound \eqref{or:pair-lower} exceeds
\eqref{or:T-pair-upper} in all nine cells.  This proves $h<90$.

\begin{table}[ht]
\centering
\caption{Closed covers for the higher $T=1$ strip.  Each cover uses
the three scale slabs $[.01,.025]$, $[.025,.045]$, and
$[.045,.061002]$.}\label{or:T-cover}
\begin{tabular}{@{}lccccc@{}}
\toprule
field assumption&moment interval&closed cells&count cells&pair cells&contours\\
\midrule
$h\ge105$&$[0,9/4]$&28&28&0&1950\\
$90\le h<105$&$[0,49/20]$&35&26&9&1881\\
\bottomrule
\end{tabular}
\end{table}

\subsection{The higher $T=1$ strip: the terminal pairs}

For $\delta=.71$ use $A=(1-\delta)^2$ when
$X\le(1-\delta)^2$, and $A=X-\kappa_\delta(X)y_0$ otherwise.
The source coefficient satisfies
\begin{equation}\label{or:Ay-terminal}
                             y_0 A\le11/20.
\end{equation}
Here is the whole-domain verification.  Put $m_0=(1-\delta)^2$.
For $X\le m_0$, use $3m_0<11/20$.  On $[m_0,.35]$ six times
the endpoint value of $\kappa_\delta/X$ is less than one, so
$y_0(X-\kappa_\delta y_0)$ increases up to $y_0=3$; certify
$11/20-3(X-3\kappa_\delta)>0$ on this full interval.
Above $.35$, the active coefficient satisfies
$\kappa_\delta\ge(5/11)X^2$, and the inactive coefficient is stronger.
Completion of the square then gives \eqref{or:Ay-terminal}.

If $|\Im S|<J_0y_0$, equation~\eqref{or:radial} implies
\[
                   \rho(N-2h)<2y_0(h+J_0).
\]
For $h_L\le h\le h_U<N_0/2$, we have $N-2h>0$, and hence
\[
 \rho hA<\frac{2y_0A\,h(h+J_0)}{N-2h}
       \le\frac{2(11/20)h_U(h_U+J_0)}{N_0-2h_U}.
\]
Use $f=(1-\rho_*^2(1.71)^2/2)^{-1}$,
$D_p=1+(\rho_*1.71/2)^2$, and
$B_U=2D_p(fh_U+1)/n_0<1$.  The resulting finite source upper is
\begin{equation}\label{or:T-terminal-upper}
 W<\frac{f\,2(11/20)h_U(h_U+J_0)/(N_0-2h_U)+\rho_*}{1-B_U}.
\end{equation}
The following two intervals cover $45<h<90$, including their seam:
\begin{equation}\label{or:T-terminal-table}
\begin{array}{c|c|c|c|c}
 \text{field interval}&\text{companion estimate}&Z&M&J_0\\ \hline
 (45,70]&\eqref{lt:Q057}&8917/1000&2109/1000&37/5\\
 {}[70,90)&\ref{lt:q05mu02}&17/2&19/10&7
\end{array}
\end{equation}
The second row applies because $q<3/70<.05$ and $\mu<1/70<.02$;
the first applies by \eqref{or:T-Q057-entry}.
In each row put $R=Z/h_L$, $m=M/h_L$.
In both rows $R>0$, $0<m<1$, and the inset gates
\eqref{or:pair-inset-gates} hold with $r=\rho_*$.
The corresponding values of $I_\Sigma$ are
\[
 I_\Sigma=
 \begin{cases}
  \dfrac{201103346909179}{112500000000000}>0,&(45,70],\\[3pt]
  \dfrac{715839963}{392000000}>0,&[70,90).
 \end{cases}
\]
Lemma~\ref{or:pair-source}, using
$\sqrt{1+.71^2}<123/100$, gives the lower
\begin{equation}\label{or:T-terminal-lower}
 \frac{1-\rho_*}{1+.71^2}+
 \frac{\{(1-\rho_*)(1-m)-\rho_*R^2/2+49/50\}^2}
 {2(1+.71^2)+4(123/100)m+2R^2}.
\end{equation}
It exceeds \eqref{or:T-terminal-upper} by more than $.6165$ in the
first row and more than $.0527$ in the second.  These rational
comparisons hold throughout the stated field intervals.
They contradict $W\ge E$ and exclude
\eqref{or:T-domain}.

\subsection{The higher strip with $X=1$: a lower field bound}

Assume now
\begin{equation}\label{or:X-domain}
             1/100\le\rho\le\rho_*,\quad X=1,\quad0<T\le1.
\end{equation}
In this case $d>1/2$, $0<y_0<2$, and
$y_0<2\Re(S/h)$ by \eqref{or:horizontal}.  We first prove
\begin{equation}\label{or:X-h50}
                                   h>50.
\end{equation}

If $T\le1/8$, apply \eqref{or:moving-field} with
$r_+=\rho_*$, $t_+=1/8$.  The exact expression $H_m$ exceeds $54$,
so $h>50$ at positive inset.  At zero inset,
$\rho=2T/(1+T^2)$ and
$\rho_*<2(1/8)/(1+(1/8)^2)$ imply $T<1/8$ strictly.
The finite homothety \eqref{or:homothety} can therefore be chosen
with $T'<1/8$.  It gives the same strict bound for the unchanged $h$.

Suppose $1/8\le T\le1$ and, for contradiction, $h\le50$.
The unphased radial identity gives
$N\rho/2<h(1+Ty_0)<50(1+2T)$, hence
\[
                           \rho<300/N\le R_*=6/125.
\]
Use $\delta=3/10$ and the polar constants
\[
 f=(1-R_*^2(1.3)^2/2)^{-1},\quad D_p=1+(R_*1.3/2)^2,
                 B_{50}=2D_p(50f+1)/n_0<1.
\]
The source has the whole-interval lower weight
\[
 b_0=\min_{t\in\{1/8,1\}}
                 \frac{t-R_*(1+t^2)/2}{t^2+.09}>0.
\]
Write $S/h=c_0+is_0$, $c_0>0$, and set
$k=\kappa_1(1,.3)=13719/34880$.  The critical comparison gives
$C<f\rho h(c_0-ky_0)$, and $c_0-ky_0>0$ follows from $C\ge0$
and its strictness.  Under $h\le50$, the denominator $1-B_{50}$
is positive, so the source inequality gives
\begin{equation}\label{or:X-Rayleigh-source}
 \rho h(c_0-ky_0)>K:=\frac{(1-B_{50})b_0-R_*}{f}
                 =\frac{1279084413230726}{1580047607421875}>0.
\end{equation}

Put $v_0=y_0/c_0\in(0,2)$, so $c_0-ky_0=c_0(1-kv_0)$.
Equation~\eqref{or:radial}, using $I(\beta)\le T\le1$, gives
\[
 \frac{N\rho}{2}<h\{c_0(R_*+v_0)+\sqrt{1-c_0^2}\}.
\]
Multiply by \eqref{or:X-Rayleigh-source}.
The product in the denominator is the Rayleigh form, at the unit
vector $(c_0,\sqrt{1-c_0^2})$, of
\[
             (1-kv_0)
             \begin{pmatrix}R_*+v_0&1/2\\1/2&0\end{pmatrix}.
\]
Its largest eigenvalue is less than $M=4/5$ for every
$0\le v_0\le2$.  Indeed the determinant of $M$ times the identity
minus this matrix is $a_2v_0^2+a_1v_0+a_0$, where
\[
 a_2=Mk-k^2/4>0,\quad
 a_1=-M(1-kR_*)+k/2,\quad a_0=M^2-MR_*-1/4,
 \qquad 4a_2a_0-a_1^2>0.
\]
This quadratic is positive on the whole real line; the other principal
minor is $M>0$.  Also $1-2k>0$.  Thus
\[
 h^2>\frac{KN_0}{2M}
      =\frac{639542206615363}{202246093750}>2500,
\]
contradicting $h\le50$.  This proves \eqref{or:X-h50}.
Consequently $q<2/h<.04$ and $\mu<2/h<.04$.
Proposition~\ref{lt:q04mu04} gives
\begin{equation}\label{or:X-low-phase}
                              |\Im S|<(15/2)y_0.
\end{equation}
The homogeneous version includes all positive $q$ in this rectangle:
the radius is $(45/4)q$, and the center bound
$Z=7.706829806638$ satisfies $Z+.45<(13/5)\pi$.

\subsection{The higher $X=1$ strip: field bounds on closed cells}

For a cell
\begin{equation}\label{or:X-cell}
       \rho\in[p,r],\qquad y_0\in[l,u],\qquad T\in[t_-,t_+],
\end{equation}
use a phase bound $J=\gamma u$, with $\gamma=4$ on a
high-field branch and $\gamma=15/2$ on the low-field branch.
For a test $\delta$, put
\[
 k=\kappa_1(t_+,\delta),\quad \overline A=1-kl,
 b_0=\omega(0,0;\delta,t_-,t_+,r),
\]
and use \eqref{or:polar-constants} with $s_+=1+\delta$.
Where $b_0>0$ the direct source gives
\begin{equation}\label{or:X-K}
 \rho h(1-ky_0)>K,\qquad
 K=\frac1f\min_{z\in\{p,r\}}
          \left\{\frac{n_0z b_0}{n_0z+2D_pb_0}-z\right\}>0.
\end{equation}
The cell conditions include positivity of $\overline A$, the polar
denominator, and $K$.  Set
\[
 G_0=\max_{l\le y_0\le u}(1-ky_0)(1+t_+y_0),\qquad
 G_1=\max_{l\le y_0\le u}y_0(1-ky_0).
\]
Each is a concave quadratic maximum, evaluated at both endpoints
and its vertex when the vertex belongs to the closed interval.
Equations~\eqref{or:X-K} and \eqref{or:radial-unphased} imply
\begin{equation}\label{or:X-source-fields}
             h>\frac{K}{r\overline A},\qquad
             h^2>\frac{KN_0}{2G_0}.
\end{equation}
For the low-field phase \eqref{or:X-low-phase}, the exact radial
identity also gives
$\rho(N-2h)<2y_0(Th+15/2)$.
Multiplying by the same positive factor $h(1-ky_0)$ gives
\[
          K(N-2h)<2y_0(1-ky_0)h(Th+15/2),\qquad h<N/2.
\]
Thus a rational $H<N_0/2$ is a valid strict lower field if
\begin{equation}\label{or:X-phase-field}
                   K(N_0-2H)\ge2G_1H(t_+H+15/2).
\end{equation}
Indeed a contrary $h\le H$ makes its left side no smaller, and its
right side no larger, than those in \eqref{or:X-phase-field}.

Besides these source bounds, every cell has the lower field bounds
\begin{equation}\label{or:X-radial-fields}
 \frac{N_0p}{2(1+t_+u)},\qquad
 \frac{N_0p/2-J}{r+t_+u},\qquad H_m
\end{equation}
from \eqref{or:radial-unphased}, \eqref{or:radial-phase}, and
\eqref{or:moving-field}.  At zero inset the last is non-strict,
which suffices for a choice $H\le H_m$.  The moving test supplies
$H_m$; the final moment calculation includes $\beta$ once.

For every inverse row choose a rational $m>0$ satisfying
\begin{equation}\label{or:X-m-cell}
              m^2\{1/4+\max(t_--1/125,0)^2\}\ge1.
\end{equation}
By \eqref{or:gap-X} this bounds $1/d$ on the whole cell.  Use
$Y=u+10^{-9}$ in \eqref{or:inverse-parameters}, and use
\eqref{or:automatic-contour}.  The rational enlargement of $u$
ensures the required strict derivative inequalities also at a closed
moment endpoint.  All rows satisfy \eqref{or:contour-gates} and
\eqref{or:odd-gate}.  Widths, slopes and depths are computed by
\eqref{or:width}, \eqref{or:counts}, and \eqref{or:depth}.
The physical nesting conditions are $R_j<R_{j+1}<2$ and increasing
slopes; running maxima of the widths remain at most the corresponding
radii.  The regions therefore have $|\Re u|\le1$ as required by
\eqref{or:weight}.

If a depth lower bound exceeds $t_+$, the cell is impossible.  This
positive depth condition forces the height budget above $4\pi$,
so Lemma~\ref{or:narrow-strip} supplies such a root.
Otherwise compute the nested source lower in \eqref{or:Abel-source}.
For an unbounded high-field cell use the correction
$\min(0,b_0-3w_1)/H$; for a low-field cell use
\begin{equation}\label{or:X-signed-correction}
 \begin{cases}
  (b_0-3w_1)/H,&b_0-3w_1<0,\\
  (b_0-3w_1)/600,&b_0-3w_1\ge0.
 \end{cases}
\end{equation}
The latter form is used only after the ceiling $h<600$ is proved.

\subsubsection{The high-field cover}

Suppose $h\ge600$.  The bounds $q<2/600<7/625$ and
$\mu<2/600<1/200$ permit
Proposition~\ref{lt:fine200}, giving $|\Im S|<4y_0$.
On each cell take $J=4u$ and $H$ below at least one of $600$ and
the three lower fields in \eqref{or:X-radial-fields}.
For $t_->0$ put
\begin{equation}\label{or:X-high-radial}
       B_r=\frac{r/2+J/N_0}{t_-(1-rt_-/2)}.
\end{equation}
A positive $U$ bounds $h/n$ if $U\ge u$, $U^2\ge B_r$, or
$Ul\ge B_r$ for $l>0$.  Apply \eqref{or:HB-cell} with
$\overline A=1-\kappa_1(t_+,\delta)l$.  The high cover has
$61$ closed cells covering $[0,2]\times[0,1]$ in $(y_0,T)$ for
the entire scale interval $[.01,.061002]$.  Each cell has either
an actual-root depth contradiction, or a strict source lower bound
greater than \eqref{or:HB-cell}.  Cells with $t_-=0$ use the
depth alternative.  It follows that $h<600$.

\subsubsection{The low-field cover}

Use \eqref{or:X-h50} and \eqref{or:X-low-phase}.
The six closed scale slabs are
\begin{equation}\label{or:X-slabs}
 [.01,.013],\ [.013,.017],\ [.017,.023],\ [.023,.033],\
 [.033,.045],\ [.045,.061002].
\end{equation}
On each cell take $J=(15/2)u$ and choose $H>J$ below at least one
of: $50$; a bound in \eqref{or:X-radial-fields}; a bound in
\eqref{or:X-source-fields}; or a value certified by
\eqref{or:X-phase-field}.  All source tests used for these choices
satisfy $1/200\le\delta\le3/4$ and
$(r(1+\delta)/2)^2<\delta$.
If $H\ge600$, the cell contradicts the preceding ceiling.
Otherwise it has either a depth contradiction or a nested
moment contradiction.  For the latter the upper denominator is
\begin{equation}\label{or:X-low-denominator}
                   B=\frac{2D_p(600f\overline A+1)}{n_0}<1.
\end{equation}
Thus every low-field source upper uses the bound $h<600$.

The six slabs have respectively $36,38,41,38,41,46$ closed cells.
Together with the $61$ high cells this makes $301$ cells.  Their
three contradiction types are $196$ moment comparisons, $57$ actual
root depth comparisons, and $48$ field comparisons.  There are
$14891$ complete inverse contours.  With the rational rounding
conventions below, the smallest moment reserve is
\begin{equation}\label{or:X-reserve}
       \frac{84739835370681964317289}
            {240500000000000000000000000}>.0003523485.
\end{equation}
Every external endpoint and every internal depth, moment, and scale
seam belongs to the closed cover.  This excludes \eqref{or:X-domain}.

\subsection{Finite verification and the entire degree ray}
\label{or:finite-verification}

We finish by specifying the finite verification behind the interval
comparisons above.  The electronic appendix contains the exact
rational parameter files and executable verifiers.  The files specify
cells, tests, radii and trims; the following procedure recomputes
their bounds.

\begin{enumerate}
\item Every integer, fraction and finite decimal is read as an exact
  rational number.  The base rows are precisely
  Table~\ref{or:base-contours}.  The higher $T$ rows are the two files
  \cert{T6251_H105_PROPOSAL.json} and
  \cert{T6251_H90_PROPOSAL.json}; each record specifies
  $(p,r,l,u,H,\delta,\overline A,U,X_+,J)$ and either the list
  $(R,\nu,W_0)$ or the pair alternative.  The higher $X$ rows are
  \cert{X6251_high_0.01_0.061002.json} and the six
  \cert{X6251_MOVING_low_} files indexed by the consecutive endpoints
  of \eqref{or:X-slabs}.  Each record specifies $(l,u,t_-,t_+,H)$,
  its contradiction type, its source tests, and its contour choices
  $(R,\nu,W_0,m)$.
\item First verify the declared parameter domains, the inverse-Cayley
  condition, every local entrance, and the asserted lower
  field using the formulae preceding the relevant cover.  For $T$
  cells verify \eqref{or:T-Xcap}, \eqref{or:T-Abar}, and the
  radial upper.  For $X$ cells verify \eqref{or:X-m-cell} and every
  source positivity and quadratic maximum used for its field.
\item For every contour recompute $\lambda,e_R,\widehat R,B_R$,
  then all of \eqref{or:contour-gates} and the appropriate horizontal
  and height inequalities.  In the higher strips use
  \eqref{or:automatic-contour}.  The displayed Taylor reciprocal is
  strictly greater than $e^{-W_0}$.  Exponential upper bounds in the
  quartic gates use, for $0\le x<32$,
  \[
    e^x\le\sum_{j=0}^{30}\frac{x^j}{j!}
           +\frac{x^{31}}{31!(1-x/32)}.
  \]
  The geometric tail proves the inequality for every omitted order.
\item In the higher covers, round slopes and weights down, and real
  widths up, to the rational lattice $10^{-12}\mathbb Z$.
  Use $7/22<1/\pi$ for slopes and $\pi<22/7$ for depth losses.
  For every $X$ row also verify $Hs_j>1$ after rounding its slope.
  Square roots needed for radial uppers or pair denominators are
  rounded up to $10^{-16}\mathbb Z$, checking the squared enclosure
  by integer arithmetic.  Take running maxima of the widths and
  future maxima of the weights, verify their nesting, and compute
  the Abel sums exactly with the signed correction
  \eqref{or:X-signed-correction} when appropriate.  The $X$ source
  upper is rounded up to $10^{-12}\mathbb Z$ before subtraction.
\item Each cell using nested counts must have $0<B<1$ and a positive
  lower-minus-upper reserve.  Each depth cell must have its rounded lower depth strictly
  greater than $t_+$.  Each low-field exclusion must have $H\ge600$.
  Each pair cell must satisfy $0<r<1$, $R>0$, $0<m<1$, $\delta>0$, $I_\Sigma>0$,
  both inset gates, both endpoint chord inequalities, its squared
  radical enclosure, and the strict
  comparison \eqref{or:pair-lower}--\eqref{or:T-pair-upper}.
\item Covering is checked as an assertion about closed rectangles.
  For each $X$ file, sort all depth endpoints.  At every endpoint
  and one point in each intervening open interval, collect the
  $[l,u]$ intervals of all cells whose closed depth interval contains
  the point.  Their sorted union must cover $[0,2]$ without a gap.
  Membership is constant between adjacent depth endpoints, so this
  proves coverage of the whole rectangle.  Check exact adjacency
  and the exterior endpoints of the six scale slabs.  For each
  $T$ cover, sort all moment endpoints and check that the active
  scale intervals cover $[.01,.061002]$ on every intervening open
  fiber; the closed cells include their endpoint fibers as well.
\end{enumerate}

Here is a finite positivity method for all rational functions used in
the $T$ entrances, source coefficients and terminal bound.  Reduce the
function to a polynomial numerator and denominator over $\mathbb Q$.
Verify nonzero endpoint values, use the Euclidean signed-remainder
sequence (Sturm's rule) to show that neither has a root in the closed
interval, and check that their signs agree at one endpoint.  At an
explicit vanishing factor such as $X^2$ or $1-X$, remove that factor
symbolically and check its endpoint separately.  This proves every
one-variable positivity statement in
\eqref{or:T-first-square}--\eqref{or:T-q-low-poly},
\eqref{or:T-Abar}, and \eqref{or:Ay-terminal} on its specified
interval.  The base source-compensation polynomial is instead
expanded in the Bernstein basis of its specified interval; all its
coefficients are positive.  To make this test explicit, if
$P(a+(b-a)t)=\sum_{j=0}^d c_jt^j$, its Bernstein coefficients are
\[
                   b_k=\sum_{j=0}^k c_j
                                   \frac{\binom{k}{j}}{\binom{d}{j}}.
\]
Nonnegative basis functions whose sum is one then prove positivity
throughout $[a,b]$.

For the two rectangle inequalities
\eqref{or:T-D51-poly}--\eqref{or:T-h45-poly}, collect the exact
coefficients as $a(X)y_0^2+b(X)y_0+c(X)$.  Subdivide the closed
$X$ interval dyadically.  On each leaf the same one-variable procedure
verifies the positive values at $y_0=0,3$ and one of the following
sufficient conditions: $b>0$; $6a+b<0$; $a<0$; or $a>0$ and
$4ac-b^2>0$.  The first condition suffices because if $a\ge0$ the
quadratic is increasing, whereas if $a<0$ its two positive endpoint
values suffice by concavity.  The second has the analogous decreasing
or concave justification.  The third and fourth follow from concavity
and completion of the square, respectively.
The coefficients $a,b$ have a factor $X$; checking their quotients
and $c(0)>0$ handles the entire $X=0$ edge.  Both children of each
subdivision are retained.  The resulting sorted leaves have exact
adjacent endpoints $0$ and $1$, proving the entire rectangle.
This proves the two coupled source and matrix inequalities over their
full rectangles.

These procedures verify the $55$ base contours and all base
inequalities; the higher $T$ procedure verifies $63$ closed cells,
$3831$ contours and $9$ pair cells; the higher $X$ procedure verifies
the $301$ closed cells and reserve \eqref{or:X-reserve}.
The least higher-$T$ reserve among cells using nested counts exceeds
$.0001411$.

Finally every argument extends from $N_0=6250$ to the entire ray
$N\ge N_0$.  All local companion, inverse-log and crescent estimates
are independent of the degree.  The direct source in
\eqref{or:direct-K} increases with $n$; radial lower fields increase
with $N$; the terms $J/N$ and $1/n$ in upper bounds decrease.
The coupled inequality \eqref{or:X-phase-field} strengthens with
$N$, and the two terminal uppers decrease with $N$.
The fixed ceilings and closed parameter covers therefore
work for every $n\ge6251$.  The base range and the two higher strips
exhaust \eqref{or:normalization}, proving
Proposition~\ref{or:theorem}.
